% !TEX program = xelatex
\documentclass[a4paper,12pt]{report}

\usepackage{fontspec}
\usepackage{xeCJK}
\IfFontExistsTF{Microsoft JhengHei}{\setCJKmainfont{Microsoft JhengHei}}{\setCJKmainfont{Noto Sans CJK TC}}
\IfFontExistsTF{Microsoft JhengHei}{\setCJKsansfont{Microsoft JhengHei}}{\setCJKsansfont{Noto Sans CJK TC}}
\IfFontExistsTF{Microsoft JhengHei}{\setCJKmonofont{Microsoft JhengHei}}{\setCJKmonofont{Noto Sans CJK TC}}
\setmainfont{Times New Roman}
\XeTeXlinebreaklocale "zh"
\XeTeXlinebreakskip = 0pt plus 1pt

\usepackage[a4paper,margin=2.5cm]{geometry}
\usepackage{array}
\usepackage{booktabs}
\usepackage{longtable}
\usepackage{enumitem}
\usepackage{graphicx}
\usepackage{xcolor}
\usepackage{listings}
\usepackage{caption}
\usepackage{fancyhdr}
\usepackage{hyperref}
\usepackage{titlesec}
\usepackage{float}

\definecolor{linegreen}{HTML}{06C755}
\definecolor{codebackground}{HTML}{F4F6F8}
\definecolor{codekeyword}{HTML}{005CC5}
\definecolor{codestring}{HTML}{A31515}
\definecolor{codecomment}{HTML}{687078}

\hypersetup{
    colorlinks=true,
    linkcolor=linegreen,
    urlcolor=blue,
    pdftitle={LINE Messaging API 註冊與應用教學文件},
    pdfauthor={413630715 王品翔}
}

\lstdefinestyle{pythonstyle}{
    language=Python,
    backgroundcolor=\color{codebackground},
    basicstyle=\ttfamily\small,
    keywordstyle=\color{codekeyword}\bfseries,
    stringstyle=\color{codestring},
    commentstyle=\color{codecomment}\itshape,
    numbers=left,
    numberstyle=\tiny\color{gray},
    numbersep=8pt,
    frame=single,
    rulecolor=\color{gray!45},
    breaklines=true,
    showstringspaces=false,
    tabsize=4,
    columns=fullflexible,
    keepspaces=true
}
\lstset{style=pythonstyle}

\renewcommand{\contentsname}{目錄}
\renewcommand{\listfigurename}{圖目錄}
\renewcommand{\listtablename}{表目錄}
\renewcommand{\lstlistlistingname}{程式碼目錄}
\renewcommand{\figurename}{圖}
\renewcommand{\tablename}{表}
\renewcommand{\chaptername}{第}
\renewcommand{\thechapter}{\arabic{chapter}}
\renewcommand{\thesection}{\arabic{chapter}.\arabic{section}}
\renewcommand{\thesubsection}{\thesection.\arabic{subsection}}
\renewcommand{\lstlistingname}{程式碼}
\setcounter{tocdepth}{2}
\setcounter{secnumdepth}{3}

\pagestyle{fancy}
\fancyhf{}
\fancyhead[L]{LINE Messaging API 教學文件}
\fancyhead[R]{413630715 王品翔}
\fancyfoot[C]{\thepage}
\setlength{\headheight}{15pt}

\titleformat{\chapter}[display]
  {\bfseries\Huge\centering}
  {第 \thechapter 章}
  {0.5em}
  {}

\title{LINE Messaging API\\註冊與應用教學文件}
\author{學號：413630715\\姓名：王品翔}
\date{2026 年 11 月 13 日}

\begin{document}

\begin{titlepage}
    \centering
    \vspace*{2cm}
    {\Huge\bfseries 專業排版與設計 期中作業\par}
    \vspace{1.5cm}
    {\LARGE\bfseries LINE Messaging API\par}
    \vspace{0.4cm}
    {\LARGE 註冊與應用教學文件\par}
    \vfill
    \begin{tabular}{rl}
        座號： & 37 \\ [0.5cm]
        班級： & 資管三B \\ [0.5cm]
        學號： & 413630715 \\ [0.5cm]
        姓名： & 王品翔 \\
    \end{tabular}
    \vfill
    {\large 2026 年 11 月 13 日\par}
\end{titlepage}

\pagenumbering{roman}
\chapter*{摘要}
本文件以 LINE Messaging API 為主題，說明如何從 LINE Developers 帳號申請開始，建立 Provider 與 Messaging API Channel，取得 Channel Secret 和 Channel access token，最後使用 Python Flask 建立可接收與回覆訊息的 Webhook 伺服器。文件同時整理環境變數、公開 HTTPS 網址、Webhook 驗證、除錯方法與安全注意事項，讓讀者能循序完成一個最小可運作的 Echo Bot。

\noindent\textbf{關鍵字：} LINE Messaging API、Webhook、Flask、Python、Access Token
\clearpage

\tableofcontents
\clearpage
\listoffigures
\clearpage
\listoftables
\clearpage
\lstlistoflistings
\clearpage

\pagenumbering{arabic}
\chapter{LINE Messaging API 與系統概念}

\section{LINE Messaging API 簡介}
LINE Messaging API 是 LINE 提供給開發者的訊息整合介面。開發者可以建立官方帳號的後端服務，接收使用者傳送的事件，並透過 API 回覆或主動推播訊息。它的核心價值在於把熟悉的聊天介面與自動化服務連接起來，使用者不需要另外安裝應用程式就能使用功能。

本作業實作的是 Echo Bot。使用者在 LINE 聊天室輸入文字後，LINE 平台會將事件以 HTTP POST 傳送到伺服器的 Webhook URL；Flask 解析請求並交由 SDK 驗證簽章，再以 reply token 回傳相同文字。這個例子雖然簡單，卻涵蓋正式機器人的重要基礎：認證、事件路由、回覆期限與伺服器部署。

\section{常見應用場景}
\begin{itemize}[leftmargin=2.5em]
    \item \textbf{客服與 FAQ：}將常見問題整理成關鍵字或意圖，提供即時回答。
    \item \textbf{通知與提醒：}透過 Push Message 傳送訂單狀態、課程提醒或系統告警。
    \item \textbf{活動與會員服務：}利用 Rich Menu 導向報名、查詢與會員功能。
    \item \textbf{資料查詢：}將資料庫、校務系統或其他 API 包裝成聊天操作。
\end{itemize}

\section{訊息流程與元件角色}
系統可以分成 LINE 使用者、LINE Platform、公開 Webhook 伺服器與 Messaging API 四個部分。使用者訊息先抵達 LINE Platform，平台建立事件 JSON 並送到 Webhook；伺服器驗證 \texttt{X-Line-Signature} 後，再使用 reply token 呼叫回覆 API。

\begin{figure}[H]
    \centering
    \fbox{\parbox{0.88\textwidth}{\centering
    使用者\quad $\longrightarrow$\quad LINE Platform\quad $\longrightarrow$\quad Flask Webhook\\[0.4cm]
    \hfill $\longleftarrow$ 回覆訊息 $\longleftarrow$\hfill}}
    \caption{LINE Echo Bot 的訊息往返流程}
    \label{fig:message-flow}
\end{figure}

\begin{table}[H]
    \centering
    \caption{主要元件與責任}
    \label{tab:components}
    \begin{tabular}{p{0.25\textwidth}p{0.62\textwidth}}
        \toprule
        元件 & 責任 \\
        \midrule
        LINE Platform & 管理使用者訊息、事件轉送與 Channel 認證。 \\
        Webhook URL & 提供公開 HTTPS POST 端點，接收事件資料。 \\
        Channel Secret & 用來驗證請求簽章，防止偽造事件。 \\
        Access Token & 伺服器呼叫 Messaging API 時使用的授權憑證。 \\
        Flask 應用程式 & 接收請求、交給 Handler 分派，並呼叫回覆 API。 \\
        \bottomrule
    \end{tabular}
\end{table}

\section{本文件的實作範圍}
本文件採用 Python 3、Flask、\texttt{line-bot-sdk} 第三版與 \texttt{python-dotenv}。實作範圍先聚焦在文字訊息回覆，藉此說明最重要的設定與驗證步驟；圖片訊息、Rich Menu 和資料庫可以在這個基礎上延伸。

\chapter{LINE Developers 註冊與 Channel 建立}

\section{建立 LINE Developers 帳號}
首先使用瀏覽器進入 LINE Developers Console，使用 LINE 帳號登入。第一次登入時須同意開發者服務的使用條款，並填寫顯示名稱、電子郵件等基本資料。登入成功後會進入 Console 首頁；若帳號尚未建立 Provider，畫面會提供建立 Provider 的入口。

\begin{enumerate}[leftmargin=3em]
    \item 開啟 LINE Developers Console，選擇以 LINE 帳號登入。
    \item 閱讀並同意開發者服務條款，完成必要的帳號資料。
    \item 在 Console 選擇 \textbf{Create a new provider}。
    \item 輸入 Provider 名稱，例如「413630715 LINE Bot」，確認建立。
\end{enumerate}

Provider 可以視為一個專案或服務擁有者。同一個 Provider 底下可以放置不同用途的 Channel，例如測試機器人、正式機器人或 LIFF 應用程式。命名時建議加入課程、專案或環境資訊，以免日後混淆。

\section{建立 Messaging API Channel}
進入剛建立的 Provider 後，選擇建立新的 Channel，類型選擇 \textbf{Messaging API}。依序填寫 Channel name、Channel description、Category、Subcategory 與電子郵件等欄位。送出後，Console 會建立一個可管理的 Messaging API Channel，並顯示 Basic settings 與 Messaging API 分頁。

\begin{table}[H]
    \centering
    \caption{建立 Channel 時的欄位用途}
    \label{tab:channel-fields}
    \begin{tabular}{p{0.28\textwidth}p{0.58\textwidth}}
        \toprule
        欄位 & 建議填寫方式 \\
        \midrule
        Channel name & 使用能辨識環境的名稱，例如 413630715 Echo Bot。 \\
        Description & 說明這個 Channel 的用途與測試範圍。 \\
        Category & 選擇最接近實際服務的類別。 \\
        Subcategory & 依照機器人功能選擇細分類別。 \\
        Email & 填寫可收信的聯絡信箱，供平台通知使用。 \\
        \bottomrule
    \end{tabular}
\end{table}

\section{取得 QR Code 與加入好友}
在 Messaging API 分頁可以看到 Bot basic ID、QR code 與加入好友相關設定。使用 LINE App 掃描 QR Code 後，即可將機器人加入好友。此時可以先加入好友，但在 Webhook 伺服器尚未啟動前，機器人還不能正確回覆訊息。

\section{申請流程中的注意事項}
Channel 的名稱和描述是管理資訊，不等於使用者看到的所有內容；真正顯示在聊天清單的名稱與圖片，還需要在相關帳號設定中另外調整。建立完成後應確認目前操作的 Provider 和 Channel 正確，避免把測試用的 Token 放到正式環境。

\chapter{憑證、Webhook 與本機環境設定}

\section{Channel Secret 與 Access Token}
在 Channel 的 Basic settings 可以找到 Channel secret。它是用於驗證 LINE Platform 請求來源的秘密字串，不能放在公開版本庫或直接寫進程式碼。Access token 則位於 Messaging API 分頁，通常要按下 Issue 或重新發行按鈕取得。Token 具有呼叫 API 的權限，同樣必須視為機密資料。

兩者用途不同：Channel Secret 用來檢查收到的 Webhook 請求是否被竄改；Channel access token 用來授權伺服器呼叫回覆訊息 API。若任一憑證外洩，應立即在 Console 重新發行或撤銷，並同步更新部署環境。

\section{Python 套件與環境變數}
在專案目錄建立虛擬環境並安裝套件：
\begin{lstlisting}[language=bash,caption={建立 Python 虛擬環境與安裝套件},label={lst:install}]
python -m venv .venv
.venv\\Scripts\\activate
python -m pip install flask line-bot-sdk python-dotenv
\end{lstlisting}

在同一個目錄建立 `.env`，內容如下。等號右側填入 LINE Developers Console 取得的實際值，檔案不可提交到 GitHub。
\begin{lstlisting}[caption={.env 憑證設定範例},label={lst:env}]
LINE_CHANNEL_SECRET=replace-with-channel-secret
LINE_CHANNEL_ACCESS_TOKEN=replace-with-channel-access-token
\end{lstlisting}

建議另外建立 `.gitignore`，至少排除 `.env`、`.venv/`、\texttt{\_\_pycache\_/} 與編譯產物，避免秘密資料和暫存檔被上傳。

\section{Webhook URL 與 HTTPS}
LINE Platform 必須能從網際網路連線到 Webhook URL，因此 `http://localhost:8000/callback` 不能直接填入 Console。開發階段可用 ngrok 將本機的 8000 port 暫時公開：
\begin{lstlisting}[language=bash,caption={使用 ngrok 暫時公開本機服務},label={lst:ngrok}]
python app.py
ngrok http 8000
\end{lstlisting}

ngrok 會產生一個 HTTPS 網址，例如 `https://example.ngrok-free.app`。在 Messaging API 的 Webhook settings 填入 `https://example.ngrok-free.app/callback`，儲存後按 Verify。測試期間若 ngrok 網址改變，Webhook URL 也必須同步更新。

正式環境則應使用具備固定網域和 TLS 憑證的雲端服務或伺服器。無論使用哪種部署方式，URL 都應使用 HTTPS，並只開放必要的 POST 路由。

\section{Webhook 設定檢查清單}
\begin{enumerate}[leftmargin=3em]
    \item Flask 程式已啟動，且主機的 8000 port 可以接受請求。
    \item Public URL 後方包含正確的 `/callback` 路徑。
    \item LINE Console 的 Use webhook 已開啟。
    \item Verify 測試成功，伺服器回覆 HTTP 200。
    \item LINE App 已將 Bot 加入好友，且沒有關閉自動回覆造成衝突。
\end{enumerate}

\chapter{Echo Bot 程式實作與測試}

\section{程式架構}
程式由四個步驟組成。第一步使用 \texttt{load\_dotenv()} 讀取本機環境變數；第二步建立 Flask 應用程式與 LINE SDK 的 Handler；第三步在 \texttt{/callback} 接收 POST 請求並交給 Handler 驗證簽章；第四步以裝飾器註冊文字訊息事件，建立回覆物件並呼叫 Messaging API。

\section{Webhook 端點}
\texttt{callback()} 先讀取 \texttt{X-Line-Signature} 標頭與原始請求內容。這裡必須使用 \texttt{request.get\_data(as\_text=True)} 取得未修改的 body，因為 SDK 需要以原始內容計算 HMAC 簽章。若簽章不正確，\texttt{InvalidSignatureError} 會被捕捉並回傳 HTTP 400；驗證成功後則回傳 \texttt{OK}。

這個檢查很重要。若直接相信外部傳入的 JSON，而沒有驗證簽章，任何人都可能偽造 LINE 事件，讓伺服器執行不應執行的操作。正式系統還應加入日誌、速率限制、錯誤追蹤與健康檢查。

\section{文字事件與回覆}
\texttt{@handler.add(MessageEvent, message=TextMessageContent)} 只接收文字訊息事件，因此圖片或貼圖事件不會進入這個函式。進入 \texttt{handle\_message()} 後，程式使用事件中的 \texttt{reply\_token} 建立 \texttt{ReplyMessageRequest}，並將 \texttt{event.message.text} 放進 \texttt{TextMessage}。最後呼叫 \texttt{reply\_message\_with\_http\_info()} 將訊息送回 LINE。

\begin{lstlisting}[caption={Python Echo Bot 核心程式},label={lst:app-core}]
from flask import Flask, abort, request
from dotenv import load_dotenv
from linebot.v3 import WebhookHandler
from linebot.v3.exceptions import InvalidSignatureError
from linebot.v3.messaging import (
    ApiClient, Configuration, MessagingApi,
    ReplyMessageRequest, TextMessage,
)
from linebot.v3.webhooks import MessageEvent, TextMessageContent
import os

load_dotenv()
app = Flask(__name__)
channel_secret = os.getenv("LINE_CHANNEL_SECRET")
access_token = os.getenv("LINE_CHANNEL_ACCESS_TOKEN")
configuration = Configuration(access_token=access_token)
handler = WebhookHandler(channel_secret)

@app.route("/callback", methods=["POST"])
def callback():
    signature = request.headers.get("X-Line-Signature", "")
    body = request.get_data(as_text=True)
    try:
        handler.handle(body, signature)
    except InvalidSignatureError:
        abort(400)
    return "OK"

@handler.add(MessageEvent, message=TextMessageContent)
def handle_message(event):
    with ApiClient(configuration) as api_client:
        messaging_api = MessagingApi(api_client)
        reply = ReplyMessageRequest(
            reply_token=event.reply_token,
            messages=[TextMessage(text=event.message.text)],
        )
        messaging_api.reply_message_with_http_info(reply)

if __name__ == "__main__":
    app.run(host="0.0.0.0", port=8000, debug=True)
\end{lstlisting}

\section{執行與測試}
設定 `.env` 後，在虛擬環境執行 `python app.py`。先使用瀏覽器或 `curl` 確認服務可以啟動，再透過 ngrok 建立公開網址。回到 LINE Developers Console 按 Verify，最後在 LINE App 傳送「測試」，若收到相同文字，即表示事件解析、簽章驗證與回覆呼叫都完成。

\begin{table}[H]
    \centering
    \caption{常見問題與排查方向}
    \label{tab:troubleshooting}
    \begin{tabular}{p{0.28\textwidth}p{0.58\textwidth}}
        \toprule
        現象 & 排查方向 \\
        \midrule
        Verify 失敗 & 確認 ngrok 正在執行、URL 含 `/callback`，並檢查 Flask port。 \\
        回傳 400 & Channel Secret 不一致，或請求內容被中介軟體修改。 \\
        收到事件但無回覆 & 檢查 Access Token、reply token 是否有效，以及 SDK 例外訊息。 \\
        沒有任何事件 & 開啟 Use webhook，確認 Bot 已加入好友且使用正確 Channel。 \\
        \bottomrule
    \end{tabular}
\end{table}

\chapter{本文件 LaTeX 原始碼}

本章列出本文件的主要 LaTeX 原始碼與實作程式，方便檢查文件結構、套件設定及程式碼行號。\texttt{lstinputlisting} 會直接讀取工作目錄中的來源檔案，因此報告內容與實際檔案保持一致。

\section{主文件 main.tex}
\lstinputlisting[caption={主文件 main.tex 原始碼},label={lst:main-source}]{main.tex}

\section{LINE Bot 程式 app.py}
\lstinputlisting[caption={LINE Bot app.py 完整原始碼},label={lst:app-source}]{app.py}

\section{章節檔案}
以下列出第一章的章節來源，其他章節採相同的拆分方式，讓長篇文件易於維護。
\lstinputlisting[caption={第一章 chap1.tex 原始碼},label={lst:chap1-source}]{chap1/chap1.tex}

\section{原始碼維護方式}
主檔案只負責套件、格式與章節順序，各章內容放在獨立的 `.tex` 檔案中。這種結構能讓作者在修改單一章節時降低誤刪其他內容的風險，也方便多人協作與 Git 版本比較。編譯時只要在 `main.tex` 所在目錄執行 XeLaTeX，即可重新產生目錄、圖目錄、表目錄與程式碼目錄。

\chapter*{結論}
本文件完成了 LINE Messaging API 的基本註冊、Channel 建立、憑證設定、Webhook 公開與 Python Echo Bot 實作。透過 Channel Secret 驗證請求、Access Token 呼叫回覆 API，以及環境變數管理機密資料，可以建立一個具備基本安全性的聊天機器人。後續可加入資料庫、Push Message、Rich Menu、Flex Message 與正式雲端部署，將範例擴充為完整服務。

\begin{thebibliography}{9}
\bibitem{line-docs} LINE Developers，\textit{Messaging API Documentation}，\url{https://developers.line.biz/en/docs/messaging-api/}。
\bibitem{flask-docs} Pallets Projects，\textit{Flask Documentation}，\url{https://flask.palletsprojects.com/}。
\bibitem{python-dotenv} python-dotenv，\textit{Project Documentation}，\url{https://pypi.org/project/python-dotenv/}。
\end{thebibliography}


\end{document}
